\documentclass[11pt]{article}
\usepackage[margin=0.85in]{geometry}
\usepackage{amsmath,booktabs,xurl,hyperref}
\usepackage[T1]{fontenc}
\setlength{\emergencystretch}{3em}
\title{Independent Scenario Retest of the Current Collision-Avoidance Controller}
\author{Pan Dynamics Collision Avoidance Research}
\date{October 1, 2026}
\begin{document}
\maketitle
\section{Outcome and scope}
This is a fresh rerun of the committed controller, without algorithm or
configuration edits. Tested commit: \path{9de49299a7c4f7d6b05e6a0dea48453fb7877a5a}.
All fourteen ideal given-path cruise baselines collide from initially separated
states. The current controller produces 0 collision cases and
0 interrupted cases in 10127 executed holds. Nominal recovery is
confirmed in 12 of fourteen cases. Recovered cases enter and hold the
tolerances by 14.15--46.20~s, including a five-second dwell.

\textbf{Design discrepancy.} During this retest the user reiterated that the
controller must always use one CLF. The tested implementation does not meet
that requirement: \path{controller/solvePredictiveControl.m}, lines 311--318,
uses \texttt{laneQuadratic} while collision rows are active and substitutes
\texttt{recoveryCostToGo} when \texttt{departure==0}. It also changes the anchor
and adds a third anchor-deviation solve in the latter branch. One active CLF
constraint per frame is not the same as one unchanged CLF definition. These
results characterize the committed branch-dependent implementation; they do
not validate the requested unified-CLF design. No correction of that design
was included in this test-only task.

\section{Protocol}
Seven deterministic threat fixtures are run sequentially at 8 and 15~m/s,
using 50-ms holds and primary horizons of 8 and 16 steps respectively, with
60 completion steps. The default optimization-node safety buffer is 0.05~m;
road boundaries are disabled. Target motion is known and observations are
exact, without noise or delay. Plant advancement uses independent ODE45 with
relative/absolute tolerances $10^{-11}/10^{-12}$. Rectangle gaps are measured
at 31 points per hold (1.667~ms spacing). This is sampled simulation evidence,
not a proof of continuous-time clearance or delayed real-time execution.

Each case runs until recovery or an 80-s observation cap. Recovery requires
five continuous seconds of sampled errors within 0.1~m lateral error,
$1^\circ$ heading error, 0.1~m/s longitudinal speed error, 0.05~m/s lateral
velocity error and 0.01~rad/s yaw-rate error, beginning no earlier than 8~s.
Completing 80~s alone is not recovery. Positive measured body clearance is the
collision criterion; the 5-cm node buffer is not imposed as a separate
intersample pass requirement.

Runtime is wall time around the complete controller call, including input
handling, initialization, model/constraint construction and every optimization
stage. Offline ODE replay, geometric audit, serialization and two preparatory
warm-ups per speed are excluded. Scenario first calls are included. MATLAB
R2026a ran with \texttt{-singleCompThread} on an AMD Ryzen 7 7800X3D. Another
project computation was present on the machine; these are measured runtimes
under shared load, not isolated-hardware benchmarks. Runtime was not inserted
as a plant actuation delay.

\section{Per-scenario results}
\begin{center}\small
\begin{tabular}{rlrrrrr}
\toprule
Speed & Scenario & End (s) & Min gap (m) & Recovery (s) & Median (ms) & Max (ms)\\
\midrule
8 & headOn & 18.05 & 0.6032 & 18.05 & 118.1 & 187.1 \\
8 & acceleratingHeadOn & 17.40 & 0.4297 & 17.40 & 118.0 & 175.6 \\
8 & brakingLead & 33.45 & 1.0618 & 33.45 & 97.8 & 373.1 \\
8 & crossing & 80.00 & 0.7418 & --- & 89.7 & 314.9 \\
8 & turningCrossing & 40.40 & 0.4401 & 40.40 & 109.4 & 233.4 \\
8 & curvedHeadOn & 18.00 & 0.7066 & 18.00 & 125.3 & 200.1 \\
8 & curvedCrossing & 46.20 & 0.5889 & 46.20 & 95.6 & 284.7 \\
15 & headOn & 14.50 & 0.0355 & 14.50 & 127.7 & 166.4 \\
15 & acceleratingHeadOn & 14.15 & 0.0314 & 14.15 & 128.9 & 178.3 \\
15 & brakingLead & 45.25 & 0.7095 & 45.25 & 96.5 & 571.3 \\
15 & crossing & 80.00 & 0.4091 & --- & 115.0 & 382.2 \\
15 & turningCrossing & 39.20 & 0.1281 & 39.20 & 144.9 & 234.2 \\
15 & curvedHeadOn & 15.45 & 0.8445 & 15.45 & 150.7 & 250.8 \\
15 & curvedCrossing & 44.30 & 0.1513 & 44.30 & 124.0 & 345.3 \\
\bottomrule
\end{tabular}
\end{center}
A dash means no confirmed recovery by the observation limit. Every recovery
time includes the full five-second dwell.

At 8~m/s, crossing ends with lateral error -621.03~m, body clearance 2.42~m and 0 cost-to-go-CLF frames. The observed trajectory does not return to nominal cruise within 80~s.

At 15~m/s, crossing ends with lateral error -629.01~m, body clearance 3.63~m and 0 cost-to-go-CLF frames. The observed trajectory does not return to nominal cruise within 80~s.

\section{Controller runtime and numerical outcomes}
Across all 10127 holds, the median full-call runtime is 107.575~ms,
the 95th percentile is 181.124~ms and the maximum is 571.328~ms.
5909 holds (58.3\%) exceed the user's 100-ms target. Thus the
current measured implementation does not meet that target.
The percentile uses linear interpolation between sorted observations.

The slowest frame is brakingLead at 15~m/s and simulation time 0.20~s.
Its recorded breakdown is:
\begin{center}
\begin{tabular}{lrr}
\toprule
Component & Time (ms) & Share (\%)\\
\midrule
Initialization & 0.764 & 0.1 \\
Problem formulation & 12.204 & 2.1 \\
PCBF solve & 29.467 & 5.2 \\
CLF solve & 527.032 & 92.2 \\
Anchor-deviation solve & 0.000 & 0.0 \\
Other controller work & 1.861 & 0.3 \\
\bottomrule
\end{tabular}
\end{center}
The slowest frame has solver exit flags [1, -7], in stage order.
There are 2 flow reinitializations and 362 frames with at
least one nonpositive solver exit flag. Finite returned vectors are executed
under the current controller contract; a nonpositive flag is not counted as
an empty-result interruption. Per-stage exit flags and timings are retained
in the data bundle. This retest adds no online acceptance checks.

On the cost-to-go branch, 3452 of 3452 frames have CLF slack within the
scaled numerical tolerance $10^{-5}\max(1,V)$; 3433 of 3433 consecutive
same-branch ODE-plant pairs meet the prescribed decrease within
$10^{-6}\max(1,V)$. This branch-specific observation cannot establish one
common CLF across branch changes.

\section{Verification and reproducibility}
The independent Python geometry recomputation agrees with the MATLAB minimum
clearance in every case to within $10^{-8}$~m; recovery dwell was independently
recomputed from the recorded errors. The campaign shell returned status 143 after the completion marker and all
fourteen JSON/MAT result pairs plus the warm-up export were written; its
termination cause was not established. Results and completeness were verified
from the saved artifacts, without claiming a zero process exit status.
Frozen MATLAB, scenario and native-kernel
hashes identify the tested implementation. No new controller changes were
made, and no historical unit-test count is claimed as a test run in this task.

Data, complete frame summaries, configurations, timing environment and methods:
\path{report/CONTROLLER_SCENARIO_RETEST_20261001/}.
The manifest records the original large JSON/MAT traces in the cache and the
exported report artifacts. Native dependencies, generated report PDFs, raw
large traces and unrelated working-tree changes are excluded from the project
commit. The reproduction file gives the exact campaign and analysis commands.
\end{document}
